Breast carcinogenesis is not a binary transition from normal to malignant tissue
but a progressive deviation from epigenetic homeostasis. Field cancerisation theory
posits that histologically normal tissue adjacent to a tumour harbours early epigenetic
alterations — detectable as increased methylation variability rather than directional
mean shifts — that precede overt neoplastic transformation. Their reproducible
characterisation across independent cohorts under strict statistical control remains
an open methodological challenge.

This thesis addresses whether DNA methylation alone can discriminate Normal from
Normal-Adjacent breast tissue across three independent cohorts (GSE69914, GSE225845,
GSE287331), spanning two Illumina array generations (HumanMethylation450K and EPIC),
under a fully leakage-controlled framework. A platform-aware preprocessing pipeline
addressed probe-design bias and statistical scale distortion. Exploratory analysis
confirmed that Normal--Adjacent differences are globally subtle -- mean $|\Delta\beta|$
from 0.009 to 0.031 -- yet consistently focal, with outlier burden ratios up to 13.3
and silhouette coefficients up to 0.34.

A multi-stage stability-based feature selection framework integrated empirical
variability filtering, biologically weighted resampling with directional stability
constraints, correlation-based redundancy pruning, and constrained genomic
diversification, yielding 5,000-CpG panels per cohort. Inter-dataset concordance
analysis revealed that two cohorts share a concordant drift axis supporting
bidirectional zero-shot transfer (AUC 0.776--0.794), while the third exhibits
systematic polarity inversion with direct consequences for cross-cohort generalisation.

A Mixed-Integer Linear Programme inspired by the 0--1 knapsack formulation
was introduced as the central methodological contribution, compressing each panel
to 50 loci under joint optimisation of discriminative performance and biological
diversity. Hard constraints enforce chromosomal balance, spatial non-redundancy,
gene-level diversity, and enrichment in COSMIC breast cancer genes.

For multi-cohort integration, polarity heterogeneity was resolved via absolute
deviation transformation from cohort-specific Normal references, followed by
NeuroCombat batch harmonisation. The resulting 50-CpG panel achieves AUC = 0.902
and balanced accuracy = 0.799 on the held-out test set ($n = 177$), mapping to
18 COSMIC breast cancer genes --- spanning tumour suppressors, chromatin regulators,
developmental transcription factors, and signalling components --- compared to 11
in the unconstrained baseline. These results demonstrate that biologically constrained
combinatorial optimisation produces compact, interpretable CpG panels that preserve
near-baseline predictive performance while embedding explicit prior knowledge.